CLIP-style dual encoders support open-vocabulary retrieval and zero-shot recognition, but they often score captions by object co-occurrence while underweighting relation, role, and binding edits. We study whether relation sensitivity can be added without replacing the pretrained embedding path. Relational Residual Bottlenecking (RRB) freezes OpenCLIP ViT-B/16 and learns a bounded, parser-free residual correction, $\zout=\mathrm{normalize}(\zbase+\alpha\drel)$, from multi-channel token bottlenecks on each branch. On staged ARO/SugarCrepe relation suites, the selected no-targeted strong-anchor checkpoint improves relation accuracy from the plain-adapter baseline's $0.6956$ to $0.7734$ while retaining COCO image-to-text R@1 of $0.7340$ and ImageNet top-1 of $0.6342$. Hard token bottlenecks and no-channel residual controls achieve relation gains only with much larger retrieval and zero-shot degradation. The experiments also give a useful negative result: structured relational counterfactuals matter, while the earlier targeted slot loss is not necessary. The claim is deliberately bounded. COCO retention remains seed-fragile, Winoground transfer is flat-to-negative, Flickr30k shows an image-to-text retrieval tax, and responsibility probes do not support semantic slot specialization.
